% Part: computability
% Chapter: computability-theory
% Section: no-universal-function

\documentclass[../../../include/open-logic-section]{subfiles}

\begin{document}

\olfileid{cmp}{thy}{nou}
\olsection{సార్వత్రిక గణనీయ ప్రమేయం లేదు}

పాక్షిక గణనీయ ప్రమేయాలన్నిటికీ సార్వత్రికమైన పాక్షిక గణనీయ ప్రమేయం
ఉన్నప్పటికీ, సర్వనిర్వచిత గణనీయ ప్రమేయాలన్నిటికీ సార్వత్రికమైన
సర్వనిర్వచిత గణనీయ ప్రమేయం లేదు.
\sourcecorrection{OLTECOMTHY-002}{ఆంగ్ల మూలంలోని మొదటి వాక్యం పాక్షిక గణనీయ ప్రమేయం వాటన్నిటికీ ``సర్వనిర్వచితం'' అని అర్థమయ్యే పదాన్ని పొరపాటుగా వాడింది. వెంటనే వచ్చే సిద్ధాంతానికీ ముందు నిరూపించిన సార్వత్రికతకూ అనుగుణంగా, అది పాక్షిక గణనీయ ప్రమేయాలన్నిటికీ ``సార్వత్రికం'' అని సరిచేశాం.}

\begin{thm}
\ollabel{thm:no-univ}
సార్వత్రిక గణనీయ ప్రమేయం ఏదీ లేదు. మరో విధంగా చెప్పాలంటే, $f(x)$
ఏదైనా సర్వనిర్వచిత గణనీయ ప్రమేయమైతే, ప్రతి~$x$కూ
$f(x)=\fn{Un}'(k,x)$ అయ్యే సహజ సంఖ్య~$k$ ఒకటి ఉండే ధర్మం గల ఏ
ప్రమేయం $\fn{Un}'(k,x)$ అయినా గణనీయం కాదు.
\end{thm}

\begin{proof}
ఇది సరళమైన వికర్ణీకరణ నిరూపణ. $\fn{Un}'(k,x)$ సర్వనిర్వచితమూ
గణనీయమూ అయితే,
\[
d(x)=\fn{Un}'(x,x)+1
\]
కూడా సర్వనిర్వచితమూ గణనీయమూ అవుతుంది. అయితే నిర్వచనం ప్రకారం,
$d(k)$ అనేది $\fn{Un}'(k,k)$కు సమానం కాదు. కనుక ప్రతి $k$కూ
$d(x)$, $\fn{Un}'(k,x)$ల విలువలు కనీసం ఒక~$x$ వద్ద, అంటే $x=k$
వద్ద, భిన్నంగా ఉంటాయి.
\end{proof}

\begin{explain}
పైనున్న \olref[uni]{thm:univ-comp}, సార్వత్రిక ప్రమేయాన్ని పాక్షికంగా
ఉండనివ్వడం అనే మూల్యం చెల్లిస్తే మాత్రమే ఈ వికర్ణీకరణ వాదనను
తప్పించుకోవచ్చని చూపుతుంది. అంటే, $\fn{Un}$ సర్వనిర్వచిత గణనీయ
ప్రమేయాలకు సార్వత్రికమే; కానీ అది సర్వనిర్వచితం కాదు. పాక్షిక
సందర్భంలో వికర్ణీకరణ వాదన పనిచేయదు.
\end{explain}

\begin{prob}
  \olref{thm:no-univ} నిరూపణలోని వికర్ణీకరణ వాదన పాక్షిక సందర్భంలో
  ఎందుకు పనిచేయదో అర్థం చేసుకోవడానికి, $f(x)\simeq
  \fn{Un}(x,x)+1$ అనే ప్రమేయాన్ని పరిగణించండి. అది పాక్షిక గణనీయ
  ప్రమేయమా? అయితే దానికి సూచిక~$e$ ఉంటుంది, అంటే $f(x)\simeq
  \fn{Un}(e,x)$. $f(e)$ గురించి ఏమి చెప్పగలరు?
\end{prob}


\end{document}
